\section{Common invariant limit under uniform deaths}
\label{sec:convergence}

\subsection{The lower class and the main theorem}
\label{subsec:convergence-statement}
{\color{red} This section is visually not good. Can be shortened, proof for Bernoulii condition needs not be stated as it is trivial.}
The convergence theorem applies to a larger class than~$\mathcal M_*$. Define
\begin{equation*}
\mathcal M_-
=
\cb{
 \mu:
 \mu\bb{\chi_A^*}
 \geq
 -\chi_A^*\bb{\mb 1}
 \text{ for every }A\Subset\Lambda
}.
\end{equation*}
Since~$\chi_A^*(\mb 1)=\prod_{i\in A}(1-p_i^\star)$,
\begin{equation}
\label{eq:m-minus-affine-characterization}
\mu\in\mathcal M_-
\quad\Longleftrightarrow\quad
\frac{\mu+\mu_{\mb 1}}2\in\mathcal M_*.
\end{equation}
In particular,~$\mathcal M_*\subseteq\mathcal M_-$. For a Bernoulli product
measure, the condition is coordinatewise. Put
\begin{equation*}
p_i^- = \bb{2p_i^\star-1}\vee0,
\qquad
\mb p^-=(p_i^-)_{i\in\Lambda}.
\end{equation*}
Then
\begin{equation}
\label{eq:m-minus-product-characterization}
\mu_{\mb p}\in\mathcal M_-
\quad\Longleftrightarrow\quad
\mb p\geq\mb p^-.
\end{equation}
Indeed, the condition for~$A=\cb{i}$ is
\begin{equation*}
p_i-p_i^\star\geq-\bb{1-p_i^\star},
\end{equation*}
which is equivalent to~$p_i\geq2p_i^\star-1$. Conversely, suppose that
$\mb p\geq\mb p^-$. When~$p_i^\star<1$, set
\begin{equation*}
x_i=\frac{p_i-p_i^\star}{1-p_i^\star}\in[-1,1].
\end{equation*}
Then
\begin{equation*}
\frac{\mu_{\mb p}(\chi_A^*)}{\chi_A^*(\mb 1)}
=
\prod_{i\in A}x_i
\geq-1.
\end{equation*}
If~$p_i^\star=1$, the singleton condition forces~$p_i=1$, and that coordinate
contributes zero to both sides of the defining inequality.
{\color{red} Polynomial growth property should probably be defined much earlier, along with the lattice notation. The finite range assumption is not needed. Also for a directed graph talking about a metric is misleading since distance from $i\to j$ can be different than from $j\to i$. Here~$B(i,r) = N_*^{r}(i)$, with the notation $N_*(A) = \bigcup_{i\in A}N_*(i)$ and~$N_*(i) = N(i)\cup \cb{i}$. No metric is needed. }
For this section, assume that~$\Lambda$ is equipped with a graph metric~$d$.
Write
\begin{equation*}
B(i,r)=\cb{j\in\Lambda:d(i,j)\leq r},
\qquad
B(A,r)=\bigcup_{i\in A}B(i,r).
\end{equation*}
We say that~$\Lambda$ has polynomial growth of exponent~$D$ if, for some
$K<\infty$,
\begin{equation*}
\abs{B(i,r)}\leq K(1+r)^D
\qquad
\text{for every }i\in\Lambda\text{ and }r\geq0.
\end{equation*}
We also assume finite range in this metric: there is~$\ell<\infty$ such that
$N(i)\subseteq B(i,\ell)$ for every~$i$.

\begin{theorem}[Common invariant limit]
\label{thm:common-invariant-limit}
Suppose that~$\Lambda$ has polynomial growth of exponent~$D$, the spin system
has patch positivity, and it has~$\varepsilon$-uniform deaths for some
$\varepsilon>0$. Then there is an invariant probability measure~$\pi$ such
that, for every local function~$f$, there is~$K_f<\infty$ satisfying
\begin{equation}
\label{eq:uniform-m-minus-convergence}
\sup_{\mu\in\mathcal M_-}
\abs{\mu\bb{P_tf}-\pi(f)}
\leq
K_f(1+t)^D e^{-\varepsilon t/2}
\qquad
\text{for every }t\geq0.
\end{equation}
Consequently,~$\mu P_t\Rightarrow\pi$ for every~$\mu\in\mathcal M_-$. The
monomial moments of~$\pi$ are
\begin{equation}
\label{eq:full-patch-limit-moments}
\pi(\chi_A)
=
\E_A\Cb{
 \prod_{P\in\mathcal P}C(P)
 \ind\bb{\abs{\mathcal P}<\infty}
}
\qquad
\text{for every }A\Subset\Lambda.
\end{equation}
\end{theorem}

\begin{corollary}[Uniform exponential ergodicity]
\label{cor:uniform-ergodicity}
Under the hypotheses of Theorem~\ref{thm:common-invariant-limit}, suppose also
that
{\color{red} the 1/2 times bolded 1 looks weird, can we just have a bolded 1/2?}
\begin{equation*}
\mb p^\star\leq\frac12\mb 1.
\end{equation*}
Then every probability measure belongs to~$\mathcal M_-$, the invariant measure
$\pi$ is unique, and
\begin{equation}
\label{eq:uniform-pointwise-convergence}
\sup_{\eta\in\Omega}
\abs{P_tf(\eta)-\pi(f)}
\leq
K_f(1+t)^D e^{-\varepsilon t/2}
\end{equation}
for every local~$f$. In particular, for every~$\gamma<\varepsilon/2$ there is
$C_{f,\gamma}<\infty$ such that
\begin{equation*}
\sup_{\eta\in\Omega}
\abs{P_tf(\eta)-\pi(f)}
\leq
C_{f,\gamma}e^{-\gamma t}.
\end{equation*}
\end{corollary}
{\color{red} Terrible explanation of the actual proof strategy. The main ingredient is comparison of patch contributions between less and more noisy spin system, since thats gives the decay because more contributions from noisier system have more cancellations and are smaller. This gives us decay on the event that there are successful interactions far into the future. The complement event term converges to a limit that does not depend on initial measure, which can be proven by undoing duality.  The rest, like exact shape of the limit is technicalities.}
The proof compares the finite-horizon patch weight with the all-time full-patch
weight. At a cut time~$T$, the difference is split according to three events:
the interaction cone leaves a finite region before~$T$, a successful
interaction occurs after~$T$, or no successful interaction occurs after~$T$.
Finite propagation controls the first term. A backward chain of patches with
outgoing terminal boundary accumulates the pure-death factor in the second. On
the third event, the end patches present at time~$T$ continue without another
successful interaction, and their centered slopes decay at rate at least
$\varepsilon$.

\subsection{Positive finite-horizon and limiting weights}
\label{subsec:positive-skeleton-weights}

Remove pure deaths at rate~$\varepsilon$ from every site and write
\begin{equation}
\label{eq:uniform-death-decomposition}
\Lcal
=
\Lcal^\varepsilon
+
\varepsilon\Ncal^{0,\mb 1}.
\end{equation}
Let~$P_t^\varepsilon$ and~$C^\varepsilon$ denote the semigroup and patch
contributions of~$\Lcal^\varepsilon$. The two systems have the same signed dual,
successful-interaction skeleton, and reference patch laws. By
Theorem~\ref{thm:pure-death-comparison},~$\Lcal^\varepsilon$ has patch
positivity and the same critical profile~$\mb p^\star$.

Fix~$A\Subset\Lambda$. For~$\mu\in\mathcal M_*$, define
\begin{equation}
\label{eq:finite-horizon-weight}
W_t^\mu
=
\prod_{P\in\mathcal B_t}C(P)\,
\mu\bb{
 \prod_{P\in\mathcal E_t}
 C\bb{\eta\bb{i(P)},P}
},
\end{equation}
and define~$W_t^{\varepsilon,\mu}$ by replacing each contribution by its
$\Lcal^\varepsilon$ counterpart. Define also
\begin{align*}
W
&=
\prod_{P\in\mathcal P}C(P)
\ind\bb{\abs{\mathcal P}<\infty},
\\
W^\varepsilon
&=
\prod_{P\in\mathcal P}C^\varepsilon(P)
\ind\bb{\abs{\mathcal P}<\infty}.
\end{align*}

\begin{lemma}[Positive weights]
\label{lem:positive-weights}
For every~$\mu\in\mathcal M_*$ and~$t\geq0$,
\begin{equation}
\label{eq:positive-weight-comparison}
0\leq W_t^\mu\leq W_t^{\varepsilon,\mu},
\qquad
\E_A\Cb{W_t^{\varepsilon,\mu}}
=
\mu\bb{P_t^\varepsilon\chi_A}
\leq1.
\end{equation}
Moreover,
\begin{equation}
\label{eq:limiting-weight-mass}
0\leq W\leq W^\varepsilon,
\qquad
\E_A\Cb{W^\varepsilon}\leq1.
\end{equation}
If a full patch~$P$ has outgoing terminal boundary, then
\begin{equation}
\label{eq:outgoing-terminal-exact-factor}
C(P)
=
e^{-\varepsilon\bb{e(P)-s(P)}}C^\varepsilon(P).
\end{equation}
\end{lemma}

\begin{proof}
For a realized finite-horizon skeleton, write
\begin{equation*}
I(\mathcal Q)=\cb{i(P):P\in\mathcal Q}
\qquad
\bb{\mathcal Q\subseteq\mathcal E_t}.
\end{equation*}
The centered end expansion gives the exact identity
\begin{equation}
\label{eq:finite-weight-centered-expansion}
W_t^\mu
=
\prod_{P\in\mathcal B_t}C(P)\sum_{\mathcal Q\subseteq\mathcal E_t}
\mu\bb{\chi_{I(\mathcal Q)}^*}
\prod_{P\in\mathcal Q}\kappa(P)
\prod_{P\in\mathcal E_t\setminus\mathcal Q}
C\bb{p_{i(P)}^\star,P}.
\end{equation}
The analogous formula for~$W_t^{\varepsilon,\mu}$ has every contribution and
slope replaced by its~$\varepsilon$ counterpart. For~$\mu\in\mathcal M_*$, the
centered moments in~\eqref{eq:finite-weight-centered-expansion} are nonnegative.
The patchwise comparison from Theorem~\ref{thm:pure-death-comparison} gives
\begin{align*}
0&\leq C(P)\leq C^\varepsilon(P),
&&P\in\mathcal B_t,
\\
0&\leq C\bb{p_{i(P)}^\star,P}
\leq C^\varepsilon\bb{p_{i(P)}^\star,P},
&&P\in\mathcal E_t,
\\
0&\leq\kappa(P)\leq\kappa^\varepsilon(P),
&&P\in\mathcal E_t.
\end{align*}
Thus every summand in~\eqref{eq:finite-weight-centered-expansion} is
nonnegative and is bounded by the corresponding~$\varepsilon$ summand. This
proves~$0\leq W_t^\mu\leq W_t^{\varepsilon,\mu}$. The patch representation for
$\Lcal^\varepsilon$ gives {\color{red} I think in general in the paper it would be more meaningful to write $\mu P_t (\chi_A)$ rather than~$\mu\bb{P_t \chi_A)}$, since we are making statements on how the spin system evolves a measure. These are of course equivalent, but it matters for presentation.}
\begin{equation*}
\E_A\Cb{W_t^{\varepsilon,\mu}}
=
\mu\bb{P_t^\varepsilon\chi_A},
\end{equation*}
and the right-hand side belongs to~$[0,1]$ because~$0\leq\chi_A\leq1$.

Pointwise patch comparison gives~$0\leq W\leq W^\varepsilon$. For
$u\to\infty$, let~$W_u^{\varepsilon,\mu_{\mb 1}}$ be the all-one finite-horizon
weight. If~$\abs{\mathcal P}<\infty$, then all its end patches eventually become
the final infinite full patches and
\begin{equation*}
W_u^{\varepsilon,\mu_{\mb 1}}\longrightarrow W^\varepsilon.
\end{equation*}
If~$\abs{\mathcal P}=\infty$, then~$W^\varepsilon=0$. Hence, everywhere,
\begin{equation*}
W^\varepsilon
\leq
\liminf_{u\to\infty}W_u^{\varepsilon,\mu_{\mb 1}}.
\end{equation*}
Fatou's lemma and the finite-horizon bound give
\begin{equation*}
\E_A\Cb{W^\varepsilon}
\leq
\liminf_{u\to\infty}
\E_A\Cb{W_u^{\varepsilon,\mu_{\mb 1}}}
\leq1.
\end{equation*}

Finally, suppose that~$P$ has outgoing terminal boundary. Consistency forces the
site to remain dual-active throughout~$[s(P),e(P))$: an empty-target death would
make the outgoing terminal interaction unsuccessful. Removing the pure deaths
changes the one-site potential from~$V_i$ to~$V_i+\varepsilon$ and changes no
sign or consistency normalizer. Therefore
\begin{equation*}
F(P)=e^{-\varepsilon\bb{e(P)-s(P)}}F^\varepsilon(P)
\end{equation*}
under the common consistent patch law, and taking expectation proves
\eqref{eq:outgoing-terminal-exact-factor}.
\end{proof}

\subsection{Spatial confinement and late interactions}
\label{subsec:spatial-confinement}
\label{subsec:late-interactions}
{\color{red} This comes out of nowhere, victim of proof strategy not being explained.}
For~$T<\infty$, let
\begin{equation*}
\mb{Cone}_T
=
\bigcup_{(i,u,S)\in\mathcal I_T}
\bb{\cb{i}\cup S}\setminus\cb{\infty}
\end{equation*}
be the set of sites reached by the successful-interaction skeleton by time~$T$.
For~$A\subseteq R\Subset\Lambda$, set
\begin{equation*}
E_T^R=\cb{\mb{Cone}_T\subseteq R},
\end{equation*}
and let~$P_t^{R,0}$ be the semigroup obtained by freezing~$R^c$ in state~$0$.
Write
\begin{equation*}
\rho_A(T,R)
=
\norm{\bb{P_T-P_T^{R,0}}\chi_A}_\infty.
\end{equation*}

\begin{lemma}[Spatial confinement]
\label{lem:spatial-confinement}
For every~$\mu\in\mathcal M_*$,
\begin{align}
0
&\leq
\E_A\Cb{W_t^\mu\ind\bb{(E_T^R)^c}}
\leq
\rho_A(T,R),
\label{eq:finite-horizon-confinement-bound}
\\
0
&\leq
\E_A\Cb{W\ind\bb{(E_T^R)^c}}
\leq
\rho_A(T,R).
\label{eq:limiting-confinement-bound}
\end{align}
\end{lemma}

\begin{proof}
Proposition~\ref{prop:confined-interaction-identity} gives
\begin{equation*}
\E_A\Cb{W_t^\mu\ind\bb{E_T^R}}
=
\mu\bb{P_{t-T}P_T^{R,0}\chi_A}.
\end{equation*}
Subtracting this identity from
$\E_A[W_t^\mu]=\mu(P_{t-T}P_T\chi_A)$ gives
\begin{equation*}
\E_A\Cb{W_t^\mu\ind\bb{(E_T^R)^c}}
=
\mu\bb{P_{t-T}\bb{P_T-P_T^{R,0}}\chi_A}.
\end{equation*}
The left-hand side is nonnegative by Lemma~\ref{lem:positive-weights}, while the
absolute value of the right-hand side is at most~$\rho_A(T,R)$ because
$\mu P_{t-T}$ is a probability measure. This proves
\eqref{eq:finite-horizon-confinement-bound}.

Apply the same bound with~$\mu=\mu_{\mb 1}$ and horizon~$u\geq T$. On
$\cb{\abs{\mathcal P}<\infty}$,
$W_u^{\mu_{\mb 1}}\to W$, while on the complementary event~$W=0$. Therefore
\begin{equation*}
W\ind\bb{(E_T^R)^c}
\leq
\liminf_{u\to\infty}
W_u^{\mu_{\mb 1}}\ind\bb{(E_T^R)^c}.
\end{equation*}
Fatou's lemma now gives~\eqref{eq:limiting-confinement-bound}.
\end{proof}

For~$0\leq T<t$, let~$L_{T,t}$ be the event that no successful interaction
occurs in~$(T,t]$, and let~$L_T$ be the event that no successful interaction
occurs after~$T$.

\begin{lemma}[Backward outgoing trail]
\label{lem:backward-outgoing-trail}
Let~$(i,u,S)$ be an {\color{red} fuck does that even mean ordinary? Stop injecting random adjectives in the prose!!} ordinary successful-interaction record. There is a finite
sequence of distinct {\color{red} Just call them patches, its confusing when you keep saying "full", when not also talking about truncated patches. There are patches and sometimes truncated patches and thats it.} full patches~$P_1,\ldots,P_n$ such that every~$P_k$ has
outgoing terminal boundary, {\color{red} it would help the reader if the condition $i(P_{k+1}) \in N_*(i(P_k))$ was written}
\begin{equation*}
e(P_1)=u,
\qquad
s(P_n)=0,
\qquad
e(P_{k+1})=s(P_k)\quad(1\leq k<n).
\end{equation*}
Consequently,
\begin{equation}
\label{eq:outgoing-trail-total-length}
\sum_{k=1}^n\bb{e(P_k)-s(P_k)}=u.
\end{equation}
\end{lemma}

\begin{proof}
The source~$i$ is dual-active immediately before~$u$, so the patch at~$i$ which
ends at~$u$ has outgoing terminal boundary; call it~$P_1$. Suppose that~$P_k$
has been chosen and~$s(P_k)>0$. If~$P_k$ begins with an incoming interaction,
let~$j$ be the source of that interaction. Since the interaction is successful,
$j$ is dual-active immediately before~$s(P_k)$, and the patch at~$j$ ending at
that time has outgoing terminal boundary. If~$P_k$ begins with an outgoing
interaction, then consistency with its outgoing terminal boundary forces the
initial interaction to be a birth, so its own site was dual-active immediately
before~$s(P_k)$; choose the preceding patch at the same site. In both cases the
chosen predecessor~$P_{k+1}$ satisfies~$e(P_{k+1})=s(P_k)$ and has outgoing
terminal boundary. The times decrease strictly. Local finiteness on~$[0,u]$
therefore forces the recursion to reach the deterministic interaction at time
zero after finitely many steps. The intervals concatenate, which proves
\eqref{eq:outgoing-trail-total-length}.
\end{proof}

{\color{red} I think it might be useful to state the main lemmas first without proofs but with explainers of their roles in the proof and morally why they are true, technical proofs can follow later once the birds eye view of the proof is clear to the reader.}

\begin{lemma}[Late interactions]
\label{lem:late-interactions}
For every~$\mu\in\mathcal M_*$,
\begin{align}
0
&\leq
\E_A\Cb{W_t^\mu\ind\bb{E_T^R\cap L_{T,t}^c}}
\leq
e^{-\varepsilon T},
\label{eq:finite-late-interaction-bound}
\\
0
&\leq
\E_A\Cb{W\ind\bb{E_T^R\cap L_T^c}}
\leq
e^{-\varepsilon T}.
\label{eq:limiting-late-interaction-bound}
\end{align}
\end{lemma}

\begin{proof}
On~$L_{T,t}^c$, let~$u\in(T,t]$ be the first successful-interaction time after
$T$, and let~$P_1,\ldots,P_n$ be the trail from
Lemma~\ref{lem:backward-outgoing-trail}. Every trail patch belongs to
$\mathcal B_t$. For a fixed~$\mathcal Q\subseteq\mathcal E_t$, compare the
corresponding summands in the centered expansions of~$W_t^\mu$ and
$W_t^{\varepsilon,\mu}$. Equation
\eqref{eq:outgoing-terminal-exact-factor} and
\eqref{eq:outgoing-trail-total-length} give
\begin{equation*}
\prod_{k=1}^n C(P_k)
=
e^{-\varepsilon u}
\prod_{k=1}^n C^\varepsilon(P_k)
\leq
e^{-\varepsilon T}
\prod_{k=1}^n C^\varepsilon(P_k).
\end{equation*}
Every remaining bulk factor, end constant, and end slope is bounded by its
$\varepsilon$ counterpart, and the centered moment multiplying the summand is
nonnegative. Summing over~$\mathcal Q$ gives the pointwise inequality
\begin{equation*}
W_t^\mu\ind\bb{L_{T,t}^c}
\leq
e^{-\varepsilon T}W_t^{\varepsilon,\mu}.
\end{equation*}
Taking expectations and using~\eqref{eq:positive-weight-comparison} proves
\eqref{eq:finite-late-interaction-bound}.

On~$\cb{\abs{\mathcal P}<\infty}\cap L_T^c$, the first successful interaction
after~$T$ has the same finite backward trail, now entirely contained in the full
patch family. The same factor-by-factor comparison gives
\begin{equation*}
W\ind\bb{L_T^c}
\leq
e^{-\varepsilon T}W^\varepsilon.
\end{equation*}
Equation~\eqref{eq:limiting-weight-mass} proves
\eqref{eq:limiting-late-interaction-bound}.
\end{proof}

\subsection{Relaxation with no late interactions}
\label{subsec:no-late-relaxation}

On~$L_{T,t}$, each end patch present at time~$T$ continues to time~$t$ without
another successful interaction. Denote this continuation by~$P^{\uparrow t}$;
the contribution includes the conditional probability of the continuation. The
same notation with~$t=\infty$ denotes the resulting full patch.

\begin{lemma}[No-late-interaction relaxation]
\label{lem:no-late-relaxation}
For~$A\subseteq R\Subset\Lambda$,~$\mu\in\mathcal M_*$, and~$0\leq T<t$,
\begin{equation}
\label{eq:no-late-relaxation-bound}
\abs{
 \E_A\Cb{W_t^\mu\ind\bb{E_T^R\cap L_{T,t}}}
 -
 \E_A\Cb{W\ind\bb{E_T^R\cap L_T}}
}
\leq
\bb{1+\abs R}e^{-\varepsilon\bb{t-T}}.
\end{equation}
\end{lemma}

\begin{proof}
{\color{red} I think this proof is completely overcomplicated. The structure is simple: on $E_T^R \cap L_{T,t}$ the end patches in~$\mathcal{E}_T$ are truncations of the end patches in~$\mathcal{E}_\infty$ and there is at most~$\abs{R}$ of them. As $t\to\infty$ we have $C(z, P^{\uparrow t}) \to C\bb{P^{\uparrow \infty}}$ exponentially fast with $t-T$. So we have a bounded number of factors, each converging exponentially. Telescoping identity should solve that in one line, no?}
Proposition~\ref{prop:no-late-continuation-identity} gives
\begin{align}
\CE{W_t^\mu\ind\bb{L_{T,t}}}{\mathcal G_T}
&=
\prod_{P\in\mathcal B_T}C(P)\,
\mu\bb{
 \prod_{P\in\mathcal E_T}
 C\bb{\eta\bb{i(P)},P^{\uparrow t}}
},
\label{eq:conditional-finite-no-late-weight}
\\
\CE{W\ind\bb{L_T}}{\mathcal G_T}
&=
\prod_{P\in\mathcal B_T}C(P)
\prod_{P\in\mathcal E_T}C\bb{P^{\uparrow\infty}}.
\label{eq:conditional-limiting-no-late-weight}
\end{align}
Put
\begin{equation*}
r_i=c_i^0(\vn)+c_i^1(\vn),
\qquad
q_i=\frac{c_i^0(\vn)}{r_i}.
\end{equation*}
The uniform death assumption gives~$r_i\geq\varepsilon$, and
Corollary~\ref{cor:critical-profile-empty-bound} gives~$p_i^\star\leq q_i$.

Fix~$P\in\mathcal E_T$, based at~$i$, and put~$h=t-T$. Write
\begin{align*}
a_t(P)&=C\bb{p_i^\star,P^{\uparrow t}},
&k_t(P)&=\partial_zC\bb{z,P^{\uparrow t}},
\\
a_t^\varepsilon(P)&=C^\varepsilon\bb{p_i^\star,P^{\uparrow t}},
&k_t^\varepsilon(P)&=\partial_zC^\varepsilon\bb{z,P^{\uparrow t}},
\\
b(P)&=C\bb{P^{\uparrow\infty}}.
\end{align*}
The continuation formula gives
\begin{align*}
a_t(P)
&=
C\bb{q_i+\bb{p_i^\star-q_i}e^{-r_i h},P},
&b(P)&=C\bb{q_i,P}.
\end{align*}
Together with the patchwise pure-death comparison, this gives
\begin{equation}
\label{eq:continued-factor-comparisons}
0\leq a_t(P)\leq a_t^\varepsilon(P),
\qquad
0\leq k_t(P)\leq e^{-\varepsilon h}k_t^\varepsilon(P),
\qquad
0\leq b(P)-a_t(P)\leq e^{-\varepsilon h}b(P).
\end{equation}
The slope inequality follows explicitly from
\begin{align*}
k_t(P)
&=e^{-r_i h}\kappa(P)
\\
&\leq
e^{-\varepsilon h}e^{-\bb{r_i-\varepsilon}h}
\kappa^\varepsilon(P)
=
e^{-\varepsilon h}k_t^\varepsilon(P).
\end{align*}
For the last inequality in~\eqref{eq:continued-factor-comparisons},
\begin{align*}
b(P)-a_t(P)
&=
\kappa(P)\bb{q_i-p_i^\star}e^{-r_i h}
\\
&\leq
e^{-\varepsilon h}b(P),
\end{align*}
because
\begin{equation*}
b(P)
=
C\bb{p_i^\star,P}
+
\kappa(P)\bb{q_i-p_i^\star}
\geq
\kappa(P)\bb{q_i-p_i^\star}.
\end{equation*}

For~$\mathcal Q\subseteq\mathcal E_T$, write
$I(\mathcal Q)=\cb{i(P):P\in\mathcal Q}$. Expanding the finite continuation
average gives
\begin{equation}
\label{eq:continued-centered-expansion}
\mu\bb{
 \prod_{P\in\mathcal E_T}
 C\bb{\eta\bb{i(P)},P^{\uparrow t}}
}
=
\sum_{\mathcal Q\subseteq\mathcal E_T}
\mu\bb{\chi_{I(\mathcal Q)}^*}
\prod_{P\in\mathcal Q}k_t(P)
\prod_{P\in\mathcal E_T\setminus\mathcal Q}a_t(P).
\end{equation}
All summands are nonnegative. If~$\mathcal Q\neq\varnothing$, then
\begin{equation*}
\prod_{P\in\mathcal Q}k_t(P)
\leq
e^{-\varepsilon h\abs{\mathcal Q}}
\prod_{P\in\mathcal Q}k_t^\varepsilon(P)
\leq
e^{-\varepsilon h}
\prod_{P\in\mathcal Q}k_t^\varepsilon(P).
\end{equation*}
Together with~$a_t(P)\leq a_t^\varepsilon(P)$, removing the term
$\mathcal Q=\varnothing$ from~\eqref{eq:continued-centered-expansion} gives
\begin{align}
0
&\leq
\mu\bb{
 \prod_{P\in\mathcal E_T}
 C\bb{\eta\bb{i(P)},P^{\uparrow t}}
}
-
\prod_{P\in\mathcal E_T}a_t(P)
\notag
\\
&\leq
e^{-\varepsilon h}
\sum_{\substack{\mathcal Q\subseteq\mathcal E_T\\\mathcal Q\neq\varnothing}}
\mu\bb{\chi_{I(\mathcal Q)}^*}
\prod_{P\in\mathcal Q}k_t^\varepsilon(P)
\prod_{P\in\mathcal E_T\setminus\mathcal Q}a_t^\varepsilon(P)
\notag
\\
&\leq
e^{-\varepsilon h}
\mu\bb{
 \prod_{P\in\mathcal E_T}
 C^\varepsilon\bb{\eta\bb{i(P)},P^{\uparrow t}}
}.
\label{eq:centered-continuation-bound}
\end{align}

Enumerate~$\mathcal E_T=\cb{P_1,\ldots,P_n}$ and write
$a_j=a_t(P_j)$ and~$b_j=b(P_j)$. Since~$0\leq a_j\leq b_j$, the exact
telescoping identity
\begin{equation*}
\prod_{j=1}^n b_j-\prod_{j=1}^n a_j
=
\sum_{j=1}^n
\bb{b_j-a_j}
\prod_{k<j}a_k
\prod_{k>j}b_k
\end{equation*}
and~\eqref{eq:continued-factor-comparisons} yield
\begin{equation}
\label{eq:constant-continuation-bound}
0
\leq
\prod_{P\in\mathcal E_T}b(P)
-
\prod_{P\in\mathcal E_T}a_t(P)
\leq
\abs{\mathcal E_T}e^{-\varepsilon h}
\prod_{P\in\mathcal E_T}b(P).
\end{equation}
Sending an end patch to its site is a bijection from~$\mathcal E_T$ to
$\mb{Cone}_T$. Hence~$\abs{\mathcal E_T}\leq\abs R$ on~$E_T^R$.

Writing~$B_T=\prod_{P\in\mathcal B_T}C(P)$, equations
\eqref{eq:conditional-finite-no-late-weight} and
\eqref{eq:conditional-limiting-no-late-weight} give the pointwise bound
\begin{align*}
&\abs{
 \CE{W_t^\mu\ind\bb{L_{T,t}}}{\mathcal G_T}
 -
 \CE{W\ind\bb{L_T}}{\mathcal G_T}
}
\\
&\qquad\leq
B_T\bb{
 \mu\bb{
  \prod_{P\in\mathcal E_T}
  C\bb{\eta\bb{i(P)},P^{\uparrow t}}
 }
 -
 \prod_{P\in\mathcal E_T}a_t(P)
}
\\
&\qquad\quad+
B_T\bb{
 \prod_{P\in\mathcal E_T}b(P)
 -
 \prod_{P\in\mathcal E_T}a_t(P)
}.
\end{align*}
These are exactly the errors in~\eqref{eq:centered-continuation-bound} and
\eqref{eq:constant-continuation-bound}. For the first error, use
$C(P)\leq C^\varepsilon(P)$ on~$\mathcal B_T$, multiply
\eqref{eq:centered-continuation-bound} by~$\ind(E_T^R)$ and the bulk factors,
and apply Proposition~\ref{prop:no-late-continuation-identity} to the
$\varepsilon$ system. Its expectation is at most
\begin{equation*}
e^{-\varepsilon h}
\E_A\Cb{W_t^{\varepsilon,\mu}\ind\bb{E_T^R\cap L_{T,t}}}
\leq
e^{-\varepsilon h}.
\end{equation*}
For the second error, multiply~\eqref{eq:constant-continuation-bound} by the
bulk factors and~$\ind(E_T^R)$. By
\eqref{eq:conditional-limiting-no-late-weight} and
\eqref{eq:limiting-weight-mass}, its expectation is at most
\begin{equation*}
\abs R e^{-\varepsilon h}
\E_A\Cb{W\ind\bb{E_T^R\cap L_T}}
\leq
\abs R e^{-\varepsilon h}.
\end{equation*}
Adding the two estimates proves~\eqref{eq:no-late-relaxation-bound}.
\end{proof}

\subsection{Completion and discussion of the initial-law condition}
\label{subsec:convergence-completion}
\label{subsec:ergodicity-corollaries}

\begin{proof}[Proof of Theorem~\ref{thm:common-invariant-limit}]
For~$\mu\in\mathcal M_*$, the patch representation and the definition of~$W$
give the exact decomposition
\begin{align*}
\mu\bb{P_t\chi_A}-\E_A\Cb{W}
={}&
\E_A\Cb{W_t^\mu\ind\bb{(E_T^R)^c}}
-
\E_A\Cb{W\ind\bb{(E_T^R)^c}}
\\
&+
\E_A\Cb{W_t^\mu\ind\bb{E_T^R\cap L_{T,t}^c}}
-
\E_A\Cb{W\ind\bb{E_T^R\cap L_T^c}}
\\
&+
\E_A\Cb{W_t^\mu\ind\bb{E_T^R\cap L_{T,t}}}
-
\E_A\Cb{W\ind\bb{E_T^R\cap L_T}}.
\end{align*}
The triangle inequality and Lemmas~\ref{lem:spatial-confinement},
\ref{lem:late-interactions}, and~\ref{lem:no-late-relaxation} give
\begin{equation}
\label{eq:three-term-convergence-bound}
\abs{
 \mu\bb{P_t\chi_A}
 -
 \E_A\Cb{W}
}
\leq
2\rho_A(T,R)
+
2e^{-\varepsilon T}
+
\bb{1+\abs R}e^{-\varepsilon\bb{t-T}}.
\end{equation}
The finite-propagation estimate in Appendix~\ref{app:finite-propagation}, applied
with exponential rate~$\varepsilon$, gives~$v<\infty$ such that, for
\begin{equation*}
R_T=B(A,vT),
\end{equation*}
one has
\begin{equation*}
\rho_A(T,R_T)\leq C_Ae^{-\varepsilon T},
\qquad
\abs{R_T}\leq C_A(1+T)^D.
\end{equation*}
Taking~$T=t/2$ in~\eqref{eq:three-term-convergence-bound} proves
\begin{equation}
\label{eq:m-star-monomial-convergence}
\sup_{\mu\in\mathcal M_*}
\abs{
 \mu\bb{P_t\chi_A}
 -
 \E_A\Cb{W}
}
\leq
K_A(1+t)^D e^{-\varepsilon t/2}.
\end{equation}

Every local function is a finite linear combination of monomials, so
\eqref{eq:m-star-monomial-convergence} gives a corresponding estimate for every
local~$f$. In particular,~$\mu_{\mb 1}P_t$ has at most one subsequential weak
limit: any two such limits agree on every monomial, hence on every local
function. Compactness of~$\Omega$ gives existence of a subsequential limit, so
$\mu_{\mb 1}P_t\Rightarrow\pi$ for a probability measure~$\pi$ satisfying
\eqref{eq:full-patch-limit-moments}. To prove invariance, let~$g$ be local and
$s\geq0$. The Feller property makes~$P_sg$ continuous, and therefore
\begin{align*}
\pi\bb{P_sg}
&=
\lim_{t\to\infty}\mu_{\mb 1}\bb{P_tP_sg}
\\
&=
\lim_{t\to\infty}\mu_{\mb 1}\bb{P_{t+s}g}
=
\pi(g).
\end{align*}
Thus~$\pi P_s=\pi$.

Let now~$\mu\in\mathcal M_-$ and set
\begin{equation*}
\overline\mu
=
\frac{\mu+\mu_{\mb 1}}2\in\mathcal M_*.
\end{equation*}
Since~$\mu=2\overline\mu-\mu_{\mb 1}$,
\begin{align*}
\abs{\mu\bb{P_tf}-\pi(f)}
&\leq
2\abs{\overline\mu\bb{P_tf}-\pi(f)}
+
\abs{\mu_{\mb 1}\bb{P_tf}-\pi(f)}
\\
&\leq
3K_f(1+t)^D e^{-\varepsilon t/2}.
\end{align*}
Absorbing the factor~$3$ into~$K_f$ proves
\eqref{eq:uniform-m-minus-convergence}.
\end{proof}

\begin{proof}[Proof of Corollary~\ref{cor:uniform-ergodicity}]
{\color{red} This should be one sentence: every point mass is in~$\mathcal{M}_-$, trivially.}
If~$p_i^\star\leq1/2$, then, for every~$A\Subset\Lambda$ and~$\eta\in\Omega$,
\begin{equation*}
\frac{\chi_A^*(\eta)}{\chi_A^*(\mb 1)}
=
\prod_{i\in A}
\frac{\eta(i)-p_i^\star}{1-p_i^\star}
\geq-1,
\end{equation*}
because every factor belongs to~$[-1,1]$. Thus every point mass, and hence every
probability measure, belongs to~$\mathcal M_-$. Equation
\eqref{eq:uniform-pointwise-convergence} follows from
Theorem~\ref{thm:common-invariant-limit}. If~$\rho$ is invariant, then
\begin{equation*}
\abs{\rho(f)-\pi(f)}
=
\abs{\rho\bb{P_tf}-\pi(f)}
\leq
\sup_{\eta\in\Omega}\abs{P_tf(\eta)-\pi(f)}
\longrightarrow0,
\end{equation*}
so~$\rho=\pi$. Finally, for~$\gamma<\varepsilon/2$,
\begin{equation*}
C_{f,\gamma}
=
K_f\sup_{t\geq0}(1+t)^D e^{-\bb{\varepsilon/2-\gamma}t}
<\infty
\end{equation*}
gives the stated exponential estimate.
\end{proof}
{\color{red} Then definition of~$\mathcal{M}_-$ should be extended to ones with arbitrary constant K like this. Mention that configs with finite nubmer of zeores are in~$\mathcal{M}_-$. Then this corollary is unnecessary.}
The coefficient~$1$ in the definition of~$\mathcal M_-$ is not essential. If a
probability measure~$\mu$ satisfies, for some~$K<\infty$,
\begin{equation}
\label{eq:finite-centered-deficit}
\mu\bb{\chi_A^*}
\geq
-K\chi_A^*\bb{\mb 1}
\qquad
\text{for every }A\Subset\Lambda,
\end{equation}
then
\begin{equation*}
\frac{\mu+K\mu_{\mb 1}}{1+K}
\in
\mathcal M_*.
\end{equation*}
Writing
\begin{equation*}
\mu
=
(1+K)\frac{\mu+K\mu_{\mb 1}}{1+K}
-K\mu_{\mb 1}
\end{equation*}
and applying the uniform~$\mathcal M_*$ estimate to both terms proves
convergence to~$\pi$ with the right-hand side of
\eqref{eq:uniform-m-minus-convergence} multiplied by at most~$1+2K$.

In particular, let~$\xi$ have zeroes only on~$F\Subset\Lambda$. The uniform
death assumption and Corollary~\ref{cor:critical-profile-empty-bound} imply
$p_i^\star<1$ for every~$i$. For every~$A\Subset\Lambda$,
\begin{equation*}
\frac{\chi_A^*(\xi)}{\chi_A^*(\mb 1)}
=
\prod_{i\in A\cap F}
\bb{-\frac{p_i^\star}{1-p_i^\star}},
\end{equation*}
so~\eqref{eq:finite-centered-deficit} holds with
\begin{equation*}
K
=
\prod_{i\in F}
\max\cb{1,\frac{p_i^\star}{1-p_i^\star}}.
\end{equation*}
Every configuration with finitely many zeroes therefore converges to the same
invariant measure, with a constant depending on the zero set. The restriction
to~$\mathcal M_-$, or more generally to a finite centered deficit, enters only
through the signs of the end-factor expansion. The argument does not identify a
dynamical obstruction to convergence from lower-density initial laws.
